Durante la última década, los patrones de consumo musical han cambiado de forma significativa hacia los entornos digitales. Actualmente, las audiencias descubren música mediante plataformas de streaming, consumen videoclips en YouTube y participan en dinámicas culturales vinculadas a la música a través de redes sociales como TikTok \cite{onderdijk2023concert}. Esta transformación se intensificó desde la pandemia de COVID-19, cuando los conciertos por transmisión en vivo y otras formas de participación virtual se convirtieron en alternativas relevantes para mantener la relación entre artistas y audiencias \cite{onderdijk2023concert}. En este contexto, la realidad virtual (VR por sus siglas en inglés) ha emergido como una alternativa para transformar la experiencia musical mediante espacios más sensoriales, interactivos e inmersivos \cite{rivas2025diseno}.

Los estudios recientes sobre música en VR han explorado principalmente cómo trasladar o enriquecer la experiencia del concierto mediante recursos inmersivos, como la navegación activa, el cambio de perspectiva, el diseño de múltiples vistas, el uso de avatares y las reacciones colectivas de la audiencia \cite{lo2023immersive,lee2025concert,yan2020social,kim2022multiview,lee2023viewing,cacciuttolo2021welcome}. No obstante, este enfoque mantiene como referencia central la presentación musical en vivo. Por ello, resulta relevante explorar un formato distinto: el videoclip musical en VR, donde la experiencia no se organiza alrededor del escenario o la audiencia, sino alrededor de la progresión temporal y expresiva de una canción.

Asimismo, algunas experiencias musicales interactivas y videojuegos rítmicos dependen de mecanismos de generación procedural, en los que el contenido se produce mediante reglas algorítmicas. Aunque estos mecanismos facilitan la creación de contenido dinámico, diversos autores señalan la importancia de conservar el control autoral mediante objetivos o indicadores de experiencia definidos por el diseñador \cite{yannakakis2011experience,khaled2013design,craveirinha2015designing}. Esta propuesta aporta una aproximación distinta al uso de VR en experiencias musicales, al trasladar el énfasis desde la simulación de conciertos hacia la construcción de un videoclip inmersivo guiado por eventos temporizados y decisiones autorales.

Para organizar esta secuenciación, se propone utilizar una Arquitectura Dirigida por Eventos (EDA por sus siglas en inglés) debido a su capacidad para coordinar sistemas reactivos y desacoplados mediante eventos, productores, consumidores y gestores de eventos. Investigaciones previas han explorado la visualización de flujos de eventos en VR \cite{oberhauser2023vredstream} y la coordinación de interacciones multisensoriales en entornos virtuales mediante enfoques basados en eventos \cite{edmunds2006mixed}. No obstante, aún se identifica una oportunidad para aplicar este enfoque al diseño de un videoclip musical en VR, donde los eventos respondan a una secuencia temporal alineada con la estructura narrativa de una canción.

A partir de esta oportunidad, el presente proyecto propone el desarrollo de un sistema de secuenciación inmersiva sincronizada basado en eventos temporizados, orientado a coordinar elementos visuales, interactivos y auditivos dentro de un videoclip musical en VR. Con el fin de evaluar el impacto de este sistema en la experiencia de usuario, se emplea el I-Group Presence Questionnaire (IPQ), una herramienta ampliamente usada para medir la percepción de presencia en entornos virtuales mediante dimensiones como presencia espacial, involucramiento y realismo \cite{berkman2021igroup}. En este contexto, la presencia espacial hace referencia a la sensación de estar físicamente situado en un Entorno Virtual (VE por sus siglas en inglés); el involucramiento, con la atención y concentración del usuario durante la experiencia con el VE; y el realismo percibido, a la valoración de qué tan creíble o comparable con la realidad resulta dicho entorno \cite{castro2026immersive}.
